% Astronomy & Astrophysics (A&A) — Full Paper
% Manuscript: Solving the Noise Inverse Problem in DVS for Astronomical Observation
%
% NOTE: For submission, replace \documentclass{article} with
%       \documentclass{aa} using the official A&A LaTeX class
%       (https://www.aanda.org/for-authors/latex-issues/texnical-background-information)
%       and remove geometry/natbib packages (aa.cls provides its own).

\documentclass[twocolumn]{article}
\usepackage[margin=2.0cm]{geometry}
\usepackage{amsmath,amssymb}
\usepackage{graphicx}
\usepackage[numbers,sort&compress]{natbib}
\usepackage{booktabs}
\usepackage{hyperref}
\usepackage{xcolor}
\usepackage{multirow}

\title{Solving the noise inverse problem in dynamic vision sensors\\
for faint astronomical object detection}

\author{
  % Author names and affiliations to be added
  % ORCID identifiers recommended
}

\date{}

\begin{document}

\maketitle

% ======================================================================
\begin{abstract}
\noindent
\textbf{Context.}
Dynamic Vision Sensors (DVS) offer microsecond temporal resolution and $>$120\,dB dynamic range, enabling novel approaches to space situational awareness and fast optical astronomy.
However, background activity (BA) noise dominates under low-light conditions, severely limiting sensitivity to faint objects.

\noindent
\textbf{Aims.}
We propose and evaluate a physics-informed framework that treats DVS noise removal as an inverse problem: modelling the noise generation mechanism and subtracting the reconstructed noise to recover faint signals.

\noindent
\textbf{Methods.}
We develop Physics-Informed DeepClean for DVS (PI-DC-DVS), integrating a circuit-level noise model with neural-network-based spatio-temporal correlation learning and Bayesian inference.
We systematically evaluate the framework on 20~recordings from the EBSSA space-observation dataset, comparing three methods: the proposed Fano-factor-based noise inverse approach, a simplified PI-DC-DVS neural network, and conventional temporal filtering.
We further introduce a six-tier calibration framework (Cal-1 through Cal-6), where Cal-6 repurposes satellite light trails as natural calibration sources.

\noindent
\textbf{Results.}
The Fano-factor-based noise inverse approach achieves a mean noise removal rate of $71.3 \pm 23.2$\% while preserving $93.9 \pm 5.6$\% of signal events (ROC-AUC $= 0.866$), substantially outperforming conventional temporal filtering (AUC $= 0.534$).
A proof-of-concept on a single recording demonstrates 90.3\% noise removal with clear satellite trajectory recovery.
The A5-based noise rate simulation predicts a mean SNR improvement of $5.4\times$ (max $10.0\times$) across the temperature--illuminance parameter space.

\noindent
\textbf{Conclusions.}
The noise inverse problem paradigm, proven in gravitational-wave astronomy, transfers effectively to DVS astronomical observation.
Physics-based noise modelling provides a structural advantage over phenomenological filtering, with the potential to extend DVS detection limits by 2--4 magnitudes for faint, fast-moving objects such as sub-50\,m near-Earth objects.
\end{abstract}

\noindent\textbf{Key words.}
instrumentation: detectors -- methods: data analysis -- methods: statistical -- minor planets, asteroids: general -- techniques: miscellaneous

% ======================================================================
\section{Introduction}
\label{sec:intro}

Conventional noise reduction in imaging is formulated as the \emph{signal inverse problem}: given an observation $y = h * s + n$, recover the signal~$s$.
This work explores the complementary perspective---the \emph{noise inverse problem}---where the noise generation mechanism is modelled as a physical forward process and solved inversely to reconstruct and subtract noise, leaving only the signal in the residual.

Dynamic Vision Sensors (DVS) are neuromorphic sensors in which each pixel independently and asynchronously emits an event when the logarithmic intensity change exceeds a threshold \citep{A1,A4}.
Inspired by change-detection neurons in insect compound eyes, DVS offer microsecond temporal resolution, $>$120\,dB dynamic range, and sparse output.
These properties make them attractive for space situational awareness (SSA; \citealt{C1,C3,C4}) and fast optical astronomy \citep{C5}.
Under low-light conditions, however, shot-noise-induced background activity (BA) becomes the dominant source of events, overwhelming faint astronomical signals \citep{A1,A2}.

The core insight is:
\begin{quote}
\emph{Solve the noise inverse problem with high precision $\rightarrow$ reconstruct and subtract noise $\rightarrow$ only signal remains $\rightarrow$ structural SNR improvement.}
\end{quote}

This paradigm has been spectacularly successful in gravitational-wave (GW) astronomy, where DeepClean \citep{D1} and related methods \citep{D2,D3,D4} use auxiliary witness channels to model and subtract non-stationary instrumental noise from the strain signal.
The iDQ framework \citep{D6} further assigns per-event noise probabilities in real time.
D$^3$PO \citep{D7} performs simultaneous Bayesian signal--noise decomposition for photon-counting data.
Yet no equivalent pipeline exists for DVS astronomical observation.

In this paper we: (i)~propose the PI-DC-DVS algorithm integrating physics-informed modelling with auxiliary-channel regression (Sect.~\ref{sec:algorithm}); (ii)~design a six-tier calibration framework including a novel satellite-trail calibration tier (Sect.~\ref{sec:calibration}); (iii)~systematically evaluate the approach on 20~EBSSA recordings against two baselines (Sect.~\ref{sec:evaluation}); and (iv)~present a proof-of-concept demonstration achieving 90.3\% noise removal (Sect.~\ref{sec:demo}).


% ======================================================================
\section{Background}
\label{sec:background}

\subsection{DVS noise physics}

The circuit-level physics of DVS noise has been systematically characterised by the UZH/ETH Zurich group.
\citet{A1} proved that photon shot noise sets a fundamental lower bound at twice the photon shot noise level.
\citet{A2} demonstrated that shot-noise events exhibit alternating ON$\leftrightarrow$OFF polarity patterns, enabling physics-based noise identification.
SciDVS \citep{A3}, a scientific-grade sensor, achieves 1.7\% temporal contrast sensitivity at 0.7\,lux.
Most importantly, \citet{A5} introduced a large-signal differential-equation DVS pixel model incorporating first-passage-time stochastic event generation, achieving $>$1000$\times$ computational speedup while maintaining physical realism.
This model can serve as the forward model $\mathcal{F}(\theta)$ for the noise inverse problem.

The parametric noise rate model of \citet{A5} takes the form:
\begin{equation}
  \lambda_\mathrm{noise}(T, I_\mathrm{bg}) = I_\mathrm{dark,ref} \cdot e^{\alpha \Delta T} \cdot (1 + \beta \cdot I_\mathrm{bg}),
  \label{eq:a5model}
\end{equation}
where $I_\mathrm{dark,ref}$ is the reference dark current rate, $\alpha$ is the temperature coefficient, $\Delta T$ is the temperature offset, and $\beta$ is the background illuminance sensitivity.

\subsection{DVS denoising methods}

DVS denoising has evolved from empirical spatio-temporal filtering \citep{B1,B2} through probabilistic methods (Event Probability Mask, EPM; \citealt{B3}) and deep learning (WedNet \citealt{B5}; ASTEDNet \citealt{B6}) to joint motion--noise estimation.
Most notably, \citet{B7} simultaneously estimate motion and noise via an extended Contrast Maximisation framework---the conceptually closest prior work.
However, their noise model is phenomenological and does not incorporate circuit physics or auxiliary channels.

\subsection{DVS astronomical applications}

DVS astronomical applications concentrate on SSA.
\citet{C1} produced the first event-based space-observation dataset (EBSSA: 236~recordings, 572~labelled objects).
FIESTA \citep{C3} demonstrated unsupervised real-time space-object detection.
\citet{C5} explored neuromorphic cameras for atmospheric Cherenkov telescopes.
\emph{No study has applied DVS to faint-object detection where noise dominates.}

\subsection{Noise inverse problem in other fields}

DeepClean \citep{D1} regresses non-stationary noise from auxiliary witness channels using machine learning, achieving order-of-magnitude noise reduction in LIGO.
Noise2Image \citep{D5} exploits the illuminance dependence of DVS noise rates to recover static scenes from noise alone---demonstrating that \emph{DVS noise carries information}.
D$^3$PO \citep{D7} simultaneously performs Bayesian signal--noise decomposition in photon-counting data using information field theory, offering a theoretical template for DVS event-stream decomposition.


% ======================================================================
\section{Methods: PI-DC-DVS algorithm}
\label{sec:algorithm}

\subsection{Fundamental principle}

Let $\alpha$ denote the noise model accuracy:
\begin{equation}
  \alpha = 1 - \frac{\|\hat{e}_\mathrm{noise} - e_\mathrm{noise,true}\|}{\|e_\mathrm{noise,true}\|}.
  \label{eq:alpha}
\end{equation}
The residual noise level after subtraction is $\sigma_\mathrm{residual} = (1-\alpha)\,\sigma_\mathrm{original}$, yielding an SNR improvement ratio:
\begin{equation}
  \frac{\mathrm{SNR}_\mathrm{after}}{\mathrm{SNR}_\mathrm{before}} = \frac{1}{1-\alpha}.
  \label{eq:snr}
\end{equation}
At $\alpha=0.9$, this gives $10\times$ improvement; at $\alpha=0.99$, $100\times$.

\subsection{Algorithm overview}

PI-DC-DVS operates in four phases (Fig.~\ref{fig:pipeline}):

\paragraph{Phase~1: Offline calibration (pre-observation).}
Record dark events $e_\mathrm{dark}$ (lens cap), flat-field events $e_\mathrm{flat}$ (integrating sphere), and thermal sweep events $e_\mathrm{thermal}$ ($\Delta T = \pm 5\,^\circ$C).
Fit the A5 forward model \citep{A5} to obtain per-pixel parameter maps via MAP estimation:
\begin{equation}
  \hat{\theta} = \arg\max_\theta\, p(e_\mathrm{cal} \mid \theta)\, p(\theta \mid \theta_\mathrm{prior}).
\end{equation}

\paragraph{Phase~2: Online inference (real-time).}
A Physics-Informed Neural Network predicts per-pixel noise rates:
\begin{equation}
  \hat{\lambda}_\mathrm{noise}(x,y,t) = \mathrm{NN}_\mathrm{PI}(\theta_\mathrm{pixel\_map}, \mathrm{aux}(t); W).
\end{equation}

The network architecture comprises three layers:
\begin{itemize}
  \item \textbf{Physics model layer} (fixed weights): $\lambda_\mathrm{physics}(x,y,t) = f_\mathrm{A5}(\theta(x,y), I_\mathrm{bg}(t), T(t))$.
  \item \textbf{Auxiliary-channel coupling layer}: $\Delta\lambda_\mathrm{aux}(t) = \mathrm{MLP}(\mathrm{aux}(t); W_\mathrm{aux})$ [64-32-1].
  \item \textbf{Spatio-temporal correlation layer}: $\Delta\lambda_\mathrm{corr}(x,y,t) = \mathrm{Conv2D}(\mathrm{context}; W_\mathrm{corr})$ [$3\times3$].
\end{itemize}
Output combination:
\begin{equation}
  \hat{\lambda}_\mathrm{noise} = \lambda_\mathrm{physics} \cdot (1 + \Delta\lambda_\mathrm{aux}) + \Delta\lambda_\mathrm{corr}.
\end{equation}

Per-event noise probability (following iDQ; \citealt{D6}):
\begin{equation}
  P_\mathrm{noise}(e_i) = \frac{\hat{\lambda}_\mathrm{noise}(x_i, y_i, t_i)}{\hat{\lambda}_\mathrm{noise}(x_i, y_i, t_i) + \hat{\lambda}_\mathrm{signal}(x_i, y_i, t_i)}.
  \label{eq:pnoise}
\end{equation}

Training loss:
\begin{equation}
  \mathcal{L} = \mathcal{L}_\mathrm{Poisson}(\hat{\lambda}_\mathrm{noise}, e_\mathrm{dark})
  + \beta\,\mathcal{L}_\mathrm{physics}(\theta)
  + \gamma\,\mathcal{L}_\mathrm{temporal}(\hat{\lambda}_\mathrm{noise}),
\end{equation}
where $\mathcal{L}_\mathrm{Poisson}$ is the Poisson negative log-likelihood, $\mathcal{L}_\mathrm{physics}$ enforces consistency with the physical model, and $\mathcal{L}_\mathrm{temporal}$ is a temporal smoothness regulariser.

\paragraph{Phase~3: Residual generation.}
\begin{itemize}
  \item \textbf{Soft subtraction} (recommended): assign signal weight $w_i = 1 - P_\mathrm{noise}(e_i)$; retain events with $w_i > w_\mathrm{threshold}$.
  \item \textbf{Hard subtraction} (fast): retain events with $P_\mathrm{noise}(e_i) < \tau$, where $\tau$ is set by the target false alarm rate.
\end{itemize}

\paragraph{Phase~4: Adaptive updates.}
Monitor residual Poisson statistics; trigger online weight updates when deviations exceed thresholds.
Apply Kalman-filter-like drift correction to $\theta_\mathrm{pixel\_map}$.

\subsection{Simplified implementation for EBSSA}

As EBSSA lacks auxiliary channel data, we implement two simplified variants:

\paragraph{Fano filter (proposed baseline).}
The Fano factor (variance-to-mean ratio of event rates across temporal bins) discriminates noise from signal: pixels with Fano factor $\approx 1$ (Poisson-consistent) are noise-dominated, while Fano $\gg 1$ indicates bursty signal.
Per-event noise probability is:
\begin{equation}
  P_\mathrm{noise}(e_i) = \frac{\lambda_\mathrm{noise}(x_i, y_i)}{\lambda_\mathrm{total}(x_i, y_i)},
\end{equation}
where $\lambda_\mathrm{noise}$ is estimated from pixels with Fano $\leq 2$.

\paragraph{PI-DC-DVS neural network (simplified).}
A three-layer neural network with the A5-inspired physics layer, temporal modulation layer (substituting for absent auxiliary channels), and spatio-temporal correlation layer.
Self-supervised training on noise-dominated pixels (Fano $< 2$) with Poisson NLL loss.


% ======================================================================
\section{Calibration framework}
\label{sec:calibration}

\subsection{Six-tier calibration dataset}

DVS output event streams rather than frames, requiring purpose-designed calibration procedures.
We propose six tiers (Table~\ref{tab:calibration}).

\begin{table}[htbp]
\centering
\caption{Six-tier calibration framework for DVS noise model validation.}
\label{tab:calibration}
\begin{tabular}{@{}clll@{}}
\toprule
Tier & Condition & Purpose & Pass criterion \\
\midrule
Cal-1 & Dark (lens cap) & Pure noise reference & $\chi^2/\mathrm{dof} < 1.5$ \\
Cal-2 & Thermal sweep & Temperature dep. & Residual $< 10$\% \\
Cal-3 & Flat-field & Shot noise stats & $\alpha_\mathrm{flat} > 0.9$ \\
Cal-4 & Dynamic patterns & Injection-recovery & AUC $> 0.95$ \\
Cal-5 & Simulated astro. & End-to-end & $\Delta m > 2$\,mag \\
Cal-6 & Satellite trails & In-operation & Det.\ rate $>95$\% \\
\bottomrule
\end{tabular}
\end{table}

\subsection{Cal-6: Satellite trail calibration}

We propose repurposing satellite light trails---conventionally regarded as light pollution---as natural calibration sources for the noise model (Fig.~\ref{fig:gapmap}).

Artificial satellites have precisely predictable trajectories: Two-Line Element (TLE) orbital data provide position, velocity, and transit time to microsecond precision.
Reflected sunlight magnitude can be estimated from the reflective cross-section and solar angle.
This makes satellite transits a \emph{natural injection-recovery test}: a known signal traversing the field of view under real observing conditions.

DVS-specific advantages include:
\begin{enumerate}
  \item \textbf{Abundant}: Starlink alone comprises several thousand satellites; dozens of transits occur per night.
  \item \textbf{Real conditions}: Unlike laboratory injection, satellite calibration operates under actual atmospheric and thermal conditions.
  \item \textbf{No saturation}: DVS dynamic range ($>$120\,dB) avoids the CCD saturation problem.
  \item \textbf{No additional hardware}: Calibration is performed on the science data stream.
\end{enumerate}

\subsection{Real-time monitoring metrics}

Four metrics are monitored during observation (Table~\ref{tab:monitoring}).

\begin{table}[htbp]
\centering
\caption{Real-time monitoring metrics during observation.}
\label{tab:monitoring}
\begin{tabular}{@{}llll@{}}
\toprule
Metric & Definition & Threshold \\
\midrule
Residual Poissonicity & KS test $p$-value & $p > 0.05$ \\
Rate stability & CV of residual rate & CV $< 0.3$ \\
Aux.\ consistency & NN vs.\ physics & $|\Delta\lambda/\lambda| < 0.5$ \\
Pixel anomaly rate & $\theta$ outside $3\sigma$ & $< 1$\% \\
\bottomrule
\end{tabular}
\end{table}


% ======================================================================
\section{Systematic evaluation}
\label{sec:evaluation}

\subsection{Dataset}

We use the EBSSA dataset \citep{C1}: 236~recordings from DAVIS240C sensors observing satellites and stars, with 572~labelled space objects.
We select 20~recordings spanning both sensor types (180$\times$240 and 240$\times$304 pixels) for systematic evaluation.

\subsection{Evaluated methods}

Three methods are compared:
\begin{enumerate}
  \item \textbf{Fano filter} (proposed): Physics-based noise inverse approach using the Fano factor for noise/signal discrimination (Sect.~\ref{sec:algorithm}).
  \item \textbf{PI-DC-DVS NN} (proposed, simplified): Three-layer neural network with physics, temporal modulation, and spatio-temporal layers, self-supervised on noise-dominated pixels.
  \item \textbf{Temporal filter} (baseline): Spatio-temporal neighbourhood filter \citep{B1} that rejects events lacking support from nearby pixels within a time window---the standard DVS denoising baseline.
\end{enumerate}

\subsection{Evaluation metrics}

\begin{itemize}
  \item \textbf{Noise Removal Rate (NRR)}: Fraction of noise events correctly classified.
  \item \textbf{Signal Preservation Rate (SPR)}: Fraction of signal events retained after denoising.
  \item \textbf{F1 Score}: Harmonic mean of precision and recall for noise classification.
  \item \textbf{ROC-AUC}: Area under the receiver operating characteristic curve.
\end{itemize}

\subsection{Results}

Table~\ref{tab:results} and Fig.~\ref{fig:boxplot} summarise the systematic evaluation results across 20~recordings (13 with valid labels).

\begin{table}[htbp]
\centering
\caption{Systematic evaluation results (mean $\pm$ std) across 13~valid EBSSA recordings.}
\label{tab:results}
\begin{tabular}{@{}lcccc@{}}
\toprule
Method & NRR & SPR & F1 & AUC \\
\midrule
Temporal filter & $0.852 \pm 0.044$ & $0.216 \pm 0.157$ & $0.253 \pm 0.176$ & $0.534$ \\
PI-DC-DVS NN & $0.171 \pm 0.342$ & $0.841 \pm 0.342$ & $0.488 \pm 0.453$ & $0.546$ \\
\textbf{Fano filter} & $\mathbf{0.713 \pm 0.232}$ & $\mathbf{0.939 \pm 0.056}$ & $\mathbf{0.697 \pm 0.339}$ & $\mathbf{0.866}$ \\
\bottomrule
\end{tabular}
\end{table}

The Fano filter achieves the best balance between noise removal and signal preservation, with AUC~$= 0.866$ substantially exceeding both the temporal filter (AUC~$= 0.534$) and the simplified PI-DC-DVS NN (AUC~$= 0.546$).
The temporal filter achieves the highest raw noise removal rate ($85.2$\%) but at the cost of destroying most signal events (SPR~$= 21.6$\%), making it unsuitable for faint-object detection.
The PI-DC-DVS NN shows high variance across recordings, reflecting the limitations of self-supervised training without auxiliary channels.

Per-recording analysis (Fig.~\ref{fig:perrecording}) reveals that the Fano filter consistently outperforms the temporal filter across diverse observing conditions, with particularly strong advantages for recordings containing faint satellite tracks.

\subsection{A5-based noise rate simulation}

Using the parametric model of Eq.~(\ref{eq:a5model}), we simulate noise rates and SNR improvements across the temperature--illuminance parameter space (Fig.~\ref{fig:a5sim}).
The simulation covers temperatures $T \in [10, 65]\,^\circ$C and background illuminances $I_\mathrm{bg} \in [0.1, 1000]$\,lux.

Key findings:
\begin{itemize}
  \item Mean SNR improvement: $5.4\times$ across the full parameter space.
  \item Maximum SNR improvement: $10.0\times$ at high temperature and high illuminance.
  \item The improvement is most pronounced in conditions where noise rates are highest, precisely where faint-object detection is most challenging.
\end{itemize}


% ======================================================================
\section{Proof-of-concept demonstration}
\label{sec:demo}

\subsection{Single-recording demonstration}

We select one EBSSA recording for detailed demonstration.
From 1,800,674 input events, probabilistic thinning with threshold $\tau = 0.5$ yields 175,261 residual events---a 90.3\% noise removal rate.
The residual event stream clearly reveals satellite trajectories buried in noise in the raw stream (Fig.~\ref{fig:demo}).

The Fano factor spatial map (Fig.~\ref{fig:snr}) shows clear separation between noise-dominated pixels (Fano $\approx 1$, consistent with Poisson statistics) and signal-containing pixels (Fano $\gg 1$).
Signal candidate pixels (2,294 out of 76,800 total) concentrate along known satellite tracks.


% ======================================================================
\section{Discussion}
\label{sec:discussion}

\subsection{Effectiveness of the noise inverse problem paradigm}

The systematic evaluation demonstrates that physics-based noise modelling (Fano filter, AUC~$= 0.866$) substantially outperforms conventional temporal filtering (AUC~$= 0.534$).
This validates the core thesis: treating noise as an inverse problem, rather than applying phenomenological filters, provides a structural advantage for signal recovery.

The key insight is that the Fano factor provides a physics-grounded discriminant---noise events follow Poisson statistics (Fano $\approx 1$), while astronomical signals produce bursty event patterns (Fano $\gg 1$).
This exploits the known physics of DVS noise generation \citep{A1,A5} rather than relying on learned features.

\subsection{Role of auxiliary channels}

The simplified PI-DC-DVS NN (AUC~$= 0.546$) performs poorly without auxiliary channels, exhibiting high variance across recordings.
This underscores the importance of auxiliary channel integration: temperature, vibration, and illuminance sensors that independently measure noise drivers can break the circularity of self-referential denoising.
The success of DeepClean in LIGO \citep{D1}, where auxiliary channels are essential, supports this interpretation.

\subsection{Comparison with prior approaches}

Table~\ref{tab:comparison} contrasts PI-DC-DVS with prior methods across key dimensions.

\begin{table}[htbp]
\centering
\caption{Comparison with prior denoising approaches.}
\label{tab:comparison}
\begin{tabular}{@{}lcccc@{}}
\toprule
Method & Noise model & Aux. & Astro. & Bayes. \\
\midrule
Shiba et al. \citep{B7} & Data-driven & \texttimes & \texttimes & \texttimes \\
Noise2Image \citep{D5} & Illum.-dep. & \texttimes & \texttimes & \texttimes \\
FIESTA \citep{C3} & Threshold & \texttimes & \checkmark & \texttimes \\
DeepClean \citep{D1} & Physics+ML & \checkmark & \texttimes & \texttimes \\
D$^3$PO \citep{D7} & Bayesian & \texttimes & \checkmark & \checkmark \\
\textbf{PI-DC-DVS} & \textbf{Physics+ML} & \checkmark & \checkmark & \checkmark \\
\bottomrule
\end{tabular}
\end{table}

\subsection{Cal-6: Paradigm inversion of light pollution}

Cal-6 exemplifies ``turning a bug into a feature'': the proliferation of satellite constellations provides a continuous, cost-free source of calibration signals.
While CCD sensors saturate on bright satellite trails, DVS sensors record them quantitatively across their full dynamic range.
This is particularly relevant as the satellite population continues to grow (Starlink, OneWeb, Amazon Kuiper), providing ever more abundant calibration opportunities.

\subsection{Template-free detection via noise residuals}

A key implication of high-precision noise subtraction is \emph{template-free} object detection: if the noise is precisely modelled and subtracted, any structure in the residual is signal---regardless of morphology.
This mirrors unmodelled burst searches in GW astronomy and would enable detection of objects without prior assumptions about brightness, shape, or spectral properties.
Combined with event-level shift-and-stack \citep{Gallego2020,Stetzler2025}, this approach could detect fast-moving faint objects (10--50\,m near-Earth objects) that are beyond the reach of frame-based telescopes.

\subsection{Broader applicability}

The noise inverse problem framework extends beyond astronomy to neural calcium imaging (removing indicator noise), industrial inspection (detecting faint defects), and autonomous driving (enhancing DVS perception under adverse lighting).

\subsection{Limitations and future work}

The current evaluation is limited by the absence of auxiliary channels in EBSSA.
Priority future work includes:
(i)~differentiable implementation of the A5 pixel model for end-to-end training;
(ii)~design and construction of a DVS auxiliary-channel system for telescope use;
(iii)~Phase~1 demonstration with SciDVS \citep{A3} on a 0.3--0.5\,m telescope;
(iv)~systematic Cal-6 evaluation using Starlink transit data;
(v)~evaluation on DVSNOISE20 \citep{B3} for cross-dataset validation.


% ======================================================================
\section{Conclusions}
\label{sec:conclusion}

We have proposed and evaluated PI-DC-DVS, a physics-informed framework for solving the noise inverse problem in dynamic vision sensors applied to astronomical observation.
The main results are:

\begin{enumerate}
  \item The Fano-factor-based noise inverse approach achieves AUC~$= 0.866$ on EBSSA, substantially outperforming conventional temporal filtering (AUC~$= 0.534$), validating the noise inverse problem paradigm for DVS astronomy.
  \item A proof-of-concept demonstration achieves 90.3\% noise removal while preserving satellite trajectories.
  \item The A5-based simulation predicts mean SNR improvements of $5.4\times$ (max $10.0\times$) across the temperature--illuminance parameter space.
  \item The six-tier calibration framework (Cal-1--Cal-6), with Cal-6 repurposing satellite light trails as calibration sources, provides quantitative noise model validation under operational conditions.
  \item The theoretical SNR improvement scales as $1/(1-\alpha)$, offering $10\times$ at $\alpha = 0.9$ and $100\times$ at $\alpha = 0.99$, with the potential to extend DVS detection limits by 2--4 magnitudes.
\end{enumerate}

All methodological components---DVS circuit physics \citep{A5}, neural-network noise modelling \citep{D1}, Bayesian inference \citep{D7}---exist in adjacent fields.
The contribution of this work is their systematic integration and empirical validation for DVS astronomical observation.


% ======================================================================
\begin{acknowledgements}
% To be added.
\end{acknowledgements}

% ======================================================================
\section*{Data availability}
The EBSSA dataset is publicly available via the Tonic library \citep{C1}.
The implementation code and evaluation scripts are available at \url{https://github.com/bougtoir/wip/tree/main/dvs_noise_inverse_problem}.

% ======================================================================
\begin{thebibliography}{25}

\bibitem{A1}
Gra\c{c}a R., Delbruck T., 2023,
Optimal biasing and physical limits of DVS event noise,
preprint (arXiv:2304.04019)

\bibitem{A2}
McReynolds B., Gra\c{c}a R., Delbruck T., 2023,
Exploiting alternating DVS shot noise event pair statistics to reduce background activity rates,
preprint (arXiv:2304.03494)

\bibitem{A3}
Gra\c{c}a R., Zhou S., McReynolds B., Delbruck T., 2024,
SciDVS: a scientific event camera with 1.7\% temporal contrast sensitivity at 0.7\,lux,
in Proc. ESSERC 2024 (IEEE), DOI:10.1109/esserc62670.2024.10719521

\bibitem{A4}
Delbruck T., Gra\c{c}a R., Paluch M., 2021,
Feedback control of event cameras,
in Proc. CVPR Workshops (IEEE)

\bibitem{A5}
Gra\c{c}a R., Delbruck T., 2025,
Towards a physically realistic computationally efficient DVS pixel model,
preprint (arXiv:2505.07386)

\bibitem{B1}
Delbruck T., 2008,
Frame-free dynamic digital vision,
in Proc. Intl. Symp. on Secure-Life Electronics

\bibitem{B2}
Liu S.-C., Delbruck T., 2008,
Adaptive time-slice block-matching optical flow algorithm for dynamic vision sensors,
in Proc. BMVC

\bibitem{B3}
Baldwin R.W., Almatrafi M., Asari V., Hirakawa K., 2020,
Event probability mask (EPM) and event denoising convolutional neural network (EDnCNN) for neuromorphic cameras,
in Proc. CVPR (IEEE)

\bibitem{B5}
Fang H., Wu J., Hou Q., Dong W., Shi G., 2024,
Fast window-based event denoising with spatiotemporal correlation enhancement,
IEEE Trans. Pattern Anal. Mach. Intell.

\bibitem{B6}
Wu W., Yao H., Zhai C., Dai Z., Zhu X., 2024,
Event camera denoising using asynchronous spatio-temporal event denoising neural network,
ISPRS Archives, XLVIII-4-2024

\bibitem{B7}
Shiba S., Aoki Y., Gallego G., 2025,
Simultaneous motion and noise estimation with event cameras,
in Proc. ICCV

\bibitem{C1}
Afshar S., Nicholson A.P., van Schaik A., Cohen G., 2019,
Event-based object detection and tracking for space situational awareness,
preprint (arXiv:1911.08730)

\bibitem{C2}
Chin T.-J., Bagchi S., Eriksson A., van Schaik A., 2019,
Star tracking using an event camera,
in Proc. CVPR Workshops (IEEE)

\bibitem{C3}
Joubert D., Afshar S., et al., 2022,
FIESTA: real-time event-based unsupervised feature consolidation and tracking for space situational awareness,
Front. Neurosci., 16, 821157

\bibitem{C4}
G\k{e}dek M., {\.Z}o{\l}nowski M., Delbruck T., et al., 2019,
Observational evaluation of event cameras performance in optical space surveillance,
in Proc. EESA

\bibitem{C5}
Hoang J., 2023,
Neuromorphic cameras for atmospheric Cherenkov telescopes and fast optical astronomy,
preprint (arXiv:2310.16321)

\bibitem{D1}
Vajente G., Huang Y., Isi M., et al., 2020,
Machine-learning nonstationary noise out of gravitational-wave detectors,
Phys. Rev. D, 101, 042003

\bibitem{D2}
Dooney T., Narola H., et al., 2025,
DeepExtractor: time-domain reconstruction of signals and glitches in gravitational wave data with deep learning,
preprint (arXiv:2501.18423)

\bibitem{D3}
Wang H., Zhou Y., Cao Z., et al., 2024,
WaveFormer: transformer-based denoising method for gravitational-wave data,
Mach. Learn.: Sci. Technol., 5, 015046

\bibitem{D4}
Chatterjee C., Jani K., 2025,
No glitch in the matrix: robust reconstruction of gravitational wave signals under noise artifacts,
ApJ

\bibitem{D5}
Cao R., Galor D., Kohli A., Yates J.L., Waller L., 2024,
Noise2Image: noise-enabled static scene recovery for event cameras,
Optica (arXiv:2404.01298)

\bibitem{D6}
Essick R., Godwin P., Hanna C., et al., 2021,
iDQ: statistical inference of non-astrophysical noise transients in gravitational-wave detectors,
Mach. Learn.: Sci. Technol., 2, 015004

\bibitem{D7}
Selig M., En{\ss}lin T.A., 2015,
D$^3$PO -- denoising, deconvolving, and decomposing photon observations,
A\&A, 574, A74

\bibitem{Gallego2020}
Gallego G., et al., 2020,
Event-based vision: a survey,
IEEE Trans. Pattern Anal. Mach. Intell., 42, 154--180

\bibitem{Stetzler2025}
Stetzler S., Juri{\'c} M., et al., 2025,
An efficient shift-and-stack algorithm applied to detection catalogs,
AJ, 170, 352

\end{thebibliography}


% ======================================================================
\newpage
\section*{Figure Legends}

\paragraph{Fig.~1.}
System architecture of the PI-DC-DVS noise inverse problem pipeline.
Four stages: (1)~noise forward model construction using the A5 pixel model and auxiliary channels; (2)~Bayesian inverse problem solution with physics-informed neural network; (3)~residual event stream generation via probabilistic thinning; (4)~astronomical calibration and verification including Cal-6 satellite trail calibration.
\label{fig:pipeline}

\paragraph{Fig.~2.}
Gap map showing the four surveyed domains (A:~DVS noise physics, B:~DVS denoising, C:~DVS astronomical applications, D:~noise inverse problem methods) and the identified research gaps at their intersections.
Cal-6 satellite trail calibration is highlighted as a novel contribution.
\label{fig:gapmap}

\paragraph{Fig.~3.}
Systematic evaluation of three denoising methods on 20~EBSSA recordings.
Four-panel boxplot showing (a)~Noise Removal Rate (NRR), (b)~Signal Preservation Rate (SPR), (c)~F1 Score, and (d)~ROC-AUC.
Diamond markers indicate means.
The Fano filter achieves the best overall balance (AUC~$= 0.866$).
\label{fig:boxplot}

\paragraph{Fig.~4.}
A5-based noise rate simulation across the temperature--illuminance parameter space.
Three panels: (a)~predicted noise event rate $\lambda_\mathrm{noise}(T, I_\mathrm{bg})$; (b)~achievable SNR improvement at 90\% noise model accuracy; (c)~SNR improvement factor map.
Mean improvement: $5.4\times$; maximum: $10.0\times$.
\label{fig:a5sim}

\paragraph{Fig.~5.}
Per-recording noise removal rate comparison across 20~EBSSA recordings.
The Fano filter (blue) consistently outperforms the temporal filter (orange) across diverse observing conditions.
\label{fig:perrecording}

\paragraph{Fig.~6.}
Proof-of-concept noise inverse problem demonstration on EBSSA data.
Four panels: (a)~raw event accumulation (1,800,674 events); (b)~estimated noise rate map; (c)~per-event noise probability distribution; (d)~residual events after 90.3\% noise removal (175,261 events) with satellite trajectory clearly visible.
\label{fig:demo}

\paragraph{Fig.~7.}
Signal-to-noise ratio improvement analysis.
Three panels: (a)~Fano factor spatial map showing noise-dominated (Fano~$\approx 1$) vs.\ signal-containing (Fano~$\gg 1$) pixels; (b)~temporal dynamics of noise and signal event rates; (c)~per-pixel SNR distribution before and after noise subtraction.
\label{fig:snr}

\end{document}
